\documentclass[11pt,a4paper]{article}
\usepackage[margin=2.5cm]{geometry}
\usepackage{amsmath,amssymb,amsthm,booktabs,graphicx,natbib,hyperref,setspace}
\setstretch{1.1}
\newtheorem{proposition}{Proposition}
\newtheorem{definition}{Definition}
\theoremstyle{remark}\newtheorem{remark}{Remark}
\newcommand{\pend}[1]{\textbf{[#1]}}

\title{Revealed Aspiration Levels: Inverse Interval Goal Programming for Recovering the Targets and Priorities Behind Public Decisions}
\author{Quang-Vinh Dang\\British University Vietnam, Hung Yen, Vietnam\\\texttt{vinh.dq4@buv.edu.vn}\quad ORCID: 0000-0002-3877-8024}
\date{Draft, \today}

\begin{document}
\maketitle

\begin{abstract}
Governments, central banks and firms announce targets, but what they pursue can only be read from what they do. We formulate the recovery of aspiration levels and goal priorities from observed sequential decisions as an inverse interval goal-programming problem. The forward model is an interval goal program with asymmetric deviation weights and an adjustment cost; a change of variables makes all optimality conditions linear in the unknowns, so the transmission from decisions to goal variables need not be known in advance. We show that optimal decisions lie at breakpoints, that decisions depend on weights only through effective weights (weight times transmission), that aspiration levels are revealed exactly when a decision stops at a band edge, and that the type of adjustment cost decides what can be learned: with a linear cost, priority weights are never point-identified, whereas with a quadratic cost (gradualism) they are point-identified under a rank condition on the observed adjustment episodes. We give an exact mixed-integer linear estimator with a strict-optimality margin and a scalable algorithm that combines parallel multistart local search with exact linear-programming and large-neighbourhood polishing. We validate the method on central banks, whose inflation targets are public: in a pre-registered test on banks not used during development, the implied target was closer to the stated target than average inflation in all three banks whose inflation had deviated from target, and across exploratory and confirmatory samples in all seven such cases; when inflation was on target, average inflation was the better estimator. Applied to aid, revealed aspiration levels range from about 0.15\% of GNI to about 1\% against a common commitment of 0.7\%, and for the five donors with verified national goals they are closer to those goals than average aid levels in four cases and equal in the fifth; applied to defence, the revealed lower edge of allies' spending rose from 1.7\% to 3.2\% of GDP after 2022; a one-goal model is not informative for corporate leverage.
\end{abstract}

\noindent\textbf{Keywords:} goal programming; inverse optimization; aspiration levels; revealed preferences; mixed-integer programming; central banks; public budgets

\section{Introduction}
Goal programming (GP) asks a decision maker for aspiration levels and priorities and returns a decision \citep{charnes1961,tamiz1998}. Many public and corporate actors announce aspiration levels (an inflation target, a defence spending guideline, an aid commitment, a leverage policy) and then decide repeatedly. The inverse question is natural and practically important: given the observed decisions, which aspiration levels and priorities would make them optimal? Comparing revealed with stated aspirations measures credibility; tracking them over time shows when priorities change.

Inverse optimization recovers parameters of an optimization model from observed solutions and is a mature field \citep{ahuja2001,chan2025review}. With noisy observations, estimators based on optimality conditions can be inconsistent, while decision-space losses are consistent but hard to compute \citep{aswani2018}. For linear programs, cost vectors are generically not point-identified \citep{ahmadi2026}. Inverse multiobjective optimization infers weights of several objectives \citep{chan2014,chanlee2018,dongzeng2021,gebken2021,carrizosa2026}. We found one paper titled inverse goal programming, which adjusts deviation weights so that a single noiseless solution becomes optimal with targets fixed \citep{saif2023}; we found no method that recovers aspiration levels and weights jointly from noisy sequential decisions. Economists infer preference parameters from policy: implicit inflation targets and loss weights of central banks \citep{dennis2006,favero2003,ozlale2003,ireland2007,shapiro2021}, welfare weights from tax schedules \citep{bourguignon2012,saez2016,hendren2020}, and inequality aversion from public service allocations \citep{behrman1987}. These approaches use smooth (usually quadratic) losses; an asymmetric, piecewise-linear loss of the kind studied by \citet{shapiro2021} is a one-goal interval goal program. A goal-programming formulation adds explicit aspiration bands, several goals with asymmetric weights, and set-valued answers when the data do not determine a point.

\paragraph{Contributions.} (i) A forward model for repeated decisions under interval goals with an adjustment cost, and a change of variables that makes the inverse problem linear apart from discrete states (Section~\ref{sec:model}). (ii) Structural results: breakpoint optimality, effective weights, exact revelation of aspiration levels at band edges, non-identifiability of weights under linear adjustment cost and a rank condition for their point identification under quadratic adjustment cost (Section~\ref{sec:ident}). (iii) An exact MILP estimator with a strict-optimality margin, and a scalable multistart algorithm with exact LP/LNS polishing (Section~\ref{sec:alg}). (iv) A pre-registered validation on central banks with public targets (Section~\ref{sec:cb}) and applications to public budgets and firms (Section~\ref{sec:apps}).

\section{Forward model}\label{sec:model}
At decision time $t$ the decision maker observes $x_t\in\mathbb{R}^K$, the values goal variables would take if the policy were unchanged, and chooses a change $D$. Goal $k$ responds linearly, $y_k(D)=x_{kt}+\beta_k D$ with $\beta_k\neq 0$. Each goal is an interval goal $[\tau_k-b_k,\tau_k+b_k]$ with weights $w_k^+$ for overshooting and $w_k^-$ for undershooting. With adjustment cost $C(D)$,
\begin{equation}\label{eq:forward}
F_t(D)=\sum_{k=1}^K\Big[w_k^+\big(y_k(D)-\tau_k-b_k\big)_+ + w_k^-\big(\tau_k-b_k-y_k(D)\big)_+\Big]+\lambda\,C(D).
\end{equation}
We study $C(D)=|D|$ (linear) and $C(D)=D^2/2$ (quadratic), and normalise $\lambda=1$ (Proposition~\ref{prop:eff}). Gradual adjustment of this kind is the behaviour of central banks that change rates in repeated small steps and of firms that adjust leverage partially towards a target \citep{flannery2006}. For $\beta_k<0$ let $\alpha_k=1/|\beta_k|$, $c_k=\alpha_k\tau_k$, $d_k=\alpha_k b_k$, $\omega_k^\pm=w_k^\pm/\alpha_k$ and $g_{kt}(D)=\alpha_k x_{kt}-D-c_k$. Then
\begin{equation}\label{eq:scaled}
F_t(D)=\sum_k\Big[\omega_k^+\big(g_{kt}(D)-d_k\big)_+ + \omega_k^-\big(-g_{kt}(D)-d_k\big)_+\Big]+\lambda\,C(D),
\end{equation}
in which, for given discrete states (which goals are above, on, inside or below their bands), the kinks are linear in $(\alpha,c,d,D)$ and the slopes are the weights. For controls that move the goal variable one-for-one (budget shares), $\alpha_k=1$ and $D$ is replaced by $-D$.

\section{Identification}\label{sec:ident}
Throughout, $w_k^\pm\ge0$, $b_k\ge0$, the adjustment-cost weight is normalised to $\lambda=1$, and every observed decision is assumed to be a strict minimiser of $F_t$ (for the quadratic cost this holds automatically; for the linear cost it is the strict-optimality margin used in estimation).

\begin{proposition}[Breakpoints]\label{prop:bp}
Let $C(D)=|D|$. Then $F_t$ is convex and piecewise linear with $F_t(D)\ge|D|$, so it attains its minimum, and every strict minimiser lies in $\mathcal{C}_t=\{0\}\cup\{D: x_{kt}+\beta_kD=\tau_k\pm b_k,\ k=1,\dots,K\}$. When all $\beta_k<0$, $\mathcal{C}_t=\{0\}\cup\{\alpha_kx_{kt}-c_k\mp d_k\}_k$.
\end{proposition}
\begin{proof}
Each term is convex and piecewise linear in $D$, and the goal terms are nonnegative, so $F_t\ge|D|$ and $F_t$ is coercive; it attains its minimum. A convex piecewise-linear function is affine between consecutive kinks; a strict minimiser cannot lie in the interior of an affine piece, so it is a kink. The kinks are $D=0$ (from $|D|$) and the points where some goal variable reaches a band edge.
\end{proof}

\begin{proposition}[Normalisation and effective weights]\label{prop:eff}
Decisions are invariant to scaling $(w,\lambda)\mapsto\kappa(w,\lambda)$, $\kappa>0$, which is why $\lambda=1$. Given the sign of each $\beta_k$, the map $(w_k,\beta_k)\mapsto(\alpha_k,\omega_k)$ with $\alpha_k=1/|\beta_k|$, $\omega_k=w_k|\beta_k|$ is a bijection, so $w_k=\alpha_k\omega_k$ is identified whenever $\alpha_k$ and $\omega_k$ are. Decisions depend on the weights only through $\omega$, the marginal goal gain per unit of policy, which is the natural object for estimation.
\end{proposition}
\begin{proof}
Both claims follow from \eqref{eq:scaled}.
\end{proof}

\begin{proposition}[Edge revelation]\label{prop:edge}
Let $\beta_k<0$. If in period $t$ the strict minimiser lies on the upper (lower) band edge of goal $k$, then $\tau_k+b_k=x_{kt}-D^*_t/\alpha_k$ (respectively $\tau_k-b_k$). Hence, if the binding goal and edge are known for two upper-edge periods of goal $k$ with distinct $x_{kt}$ and one lower-edge period, then $\alpha_k$, $\tau_k$ and $b_k$ are determined.
\end{proposition}
\begin{proof}
On the upper edge $D^*_t=\alpha_kx_{kt}-(c_k+d_k)$, a linear equation in $(\alpha_k,c_k+d_k)$; two such equations with distinct $x_{kt}$ determine both. A lower-edge period then gives $c_k-d_k=\alpha_kx_{kt}-D^*_t$, hence $c_k$, $d_k$, $\tau_k=c_k/\alpha_k$ and $b_k=d_k/\alpha_k$.
\end{proof}
The condition is sufficient, not necessary, and the binding goal and edge of a period are latent in practice; the estimator in Section~\ref{sec:alg} treats them as discrete states.

\begin{proposition}[Weights are never point-identified under linear cost]\label{prop:set}
Let $C(D)=|D|$ and fix $(\alpha,c,d)$. The set of weights $\omega\ge0$ under which every observed decision is a strict minimiser is defined by finitely many strict linear inequalities. If it is nonempty it is a nonempty open polyhedron, so it contains a neighbourhood of any of its points and the weights are not point-identified. Jointly over $(\alpha,c,d,\omega)$ the identified set is a finite union of such sets and need not be convex.
\end{proposition}
\begin{proof}
For fixed $(\alpha,c,d)$ the state of each goal at $D^*_t$ (strictly above, on the edge, strictly inside, strictly below) is fixed, and the one-sided derivatives of $F_t$ at $D^*_t$ are affine in $\omega$ with coefficients determined by these states. Strict optimality is right derivative $>0$ and left derivative $<0$ in every period: finitely many strict linear inequalities, whose solution set is open. Varying $(\alpha,c,d)$ changes the states only at finitely many hyperplanes, which gives the finite union.
\end{proof}

\begin{proposition}[Weights under gradualism]\label{prop:quad}
Let $C(D)=D^2/2$ and $\beta_k<0$ for all $k$. Then $F_t$ is $1$-strongly convex, its minimiser $D^*_t$ is unique and continuous in the parameters, and $D^*_t=\sum_{k\in A_t}\omega_k^+-\sum_{k\in B_t}\omega_k^-+\sum_k\pi_{kt}$, where $A_t$ ($B_t$) is the set of goals strictly above (below) their bands at $D^*_t$ and $\pi_{kt}$ are multipliers of goals sitting exactly on an edge. In particular, if at $D^*_t$ exactly one goal $k$ is strictly above its band and every other goal is strictly inside its band, then $\omega_k^+=D^*_t$; if exactly one goal is strictly below, $\omega_k^-=-D^*_t$. Given $(\alpha,c,d)$, the weights are point-identified if and only if the incidence matrix of the sets $(A_t,B_t)$ over the periods in which no goal sits on an edge has full column rank.
\end{proposition}
\begin{proof}
$F_t-D^2/2$ is convex, so $F_t$ is $1$-strongly convex in $D$ uniformly in the parameters, and it is Lipschitz in the parameters on bounded sets; hence the minimiser is unique and continuous. Stationarity $0\in\partial F_t(D^*_t)$ gives the expression, with $p_k=\omega_k^+$ for goals strictly above, $q_k=\omega_k^-$ for goals strictly below, zero for goals strictly inside, and a multiplier in an interval for goals on an edge. In periods without edge goals the equations $D^*_t=\sum_{A_t}\omega_k^+-\sum_{B_t}\omega_k^-$ are linear in $\omega$ with incidence coefficients $\pm1$, and a linear system has a unique solution if and only if its matrix has full column rank.
\end{proof}
\begin{remark}
Under gradualism the decision maker moves in constant steps $\omega_k^\pm$ while a single goal is outside its band and the others are inside, and stops at the band edge. With two goals outside their bands only the sum of their weights is revealed. Weights and aspiration levels together are therefore identified when the data contain both single-goal adjustment periods and edge-stopping periods. For controls that move a share one-for-one ($\alpha=1$, $D$ replaced by $-D$), $\omega^+$ is the size of a cut when the share is above its band and $\omega^-$ the size of an increase when it is below.
\end{remark}

\section{Estimation}\label{sec:alg}
Observed decisions are $D_t^{\mathrm{obs}}=D_t^*+\varepsilon_t$. We minimise the decision-space loss $\frac1T\sum_t|D_t^{\mathrm{obs}}-D^*_t(\theta)|$ over $\theta=(\alpha,c,d,\omega^+,\omega^-)$, which is consistent under standard conditions but hard in general \citep{aswani2018}.

\paragraph{Exact MILP.} Binary states per period and goal (upper active, strictly above, lower active, strictly below), big-$M$ links to the linear kink functions, and stationarity written with multipliers give a mixed-integer linear program. Products of weights and binaries are linearised exactly. Under linear cost, flat optima allow the fitted decisions to absorb noise; we therefore impose a strict-optimality margin: the right derivative of $F_t$ at $D_t$ is at least $\epsilon$ and the left derivative at most $-\epsilon$. Under quadratic cost the optimum is unique and no margin is needed.

\paragraph{Multistart with exact polishing.} Under quadratic cost the decision loss is continuous in $\theta$, and $D^*_t(\theta)$ is computed exactly by evaluating $F_t$ at the $2K$ breakpoints and at the stationary points of the pieces (our implementation enumerates the $3^K$ possible slope sums; sorting the breakpoints gives $O(K\log K)$). We run Powell's method from many random starts in parallel, then polish the best solutions: with all states fixed the estimation problem is a linear program; large-neighbourhood search frees the states of a subset of periods (largest residuals plus random periods) and solves the small MILP exactly.

\paragraph{Synthetic evidence.} Table~\ref{tab:synth} summarises runs with $K=2$, $T=80$ and Laplace decision noise. The full MILP stalls on weak relaxations; differential evolution finds the noiseless optimum but not the noisy one; the multistart algorithm reaches a loss below that of the true parameters at noise $0.05$ and recovers aspiration levels, bands, transmission and, approximately, weights.
\begin{table}[t]\centering\small
\begin{tabular}{llrr}\toprule
Model & Algorithm & Fitted loss & Loss at true $\theta$\\\midrule
Linear cost, noise 0.1 & full MILP (120 s) & 0.205 & 0.089\\
 & LNS (29 s) & 0.189 & 0.089\\
 & DE + LP + LNS (32 s) & 0.101 & 0.089\\
Quadratic cost, noise 0 & DE, 3000 generations (189 s) & 0.000 & 0.000\\
Quadratic cost, noise 0.05 & DE + LP + LNS & 0.093 & 0.055\\
 & multistart + LP + LNS (267 s) & 0.052 & 0.055\\
Quadratic cost, noise 0.1 & multistart + LP + LNS (925 s) & 0.111 & 0.106\\
\bottomrule\end{tabular}
\caption{Estimation on synthetic data ($K=2$, $T=80$).}\label{tab:synth}
\end{table}

\section{Validation on central banks}\label{sec:cb}
Inflation-targeting central banks publish numeric targets, so the implied target $\hat\tau$ can be compared with the truth. Inputs are quarterly: the change in the policy rate (BIS), inflation four quarters ahead (realised, as a proxy for the forecast) and the GDP growth gap; quarters with the policy rate at or below 0.25\% are dropped. The benchmark is average inflation over the same window, which is a strong estimator for credible banks because inflation averages close to target.

\paragraph{Development.} Earlier specifications with a linear adjustment cost did not beat the benchmark (mean absolute error 0.72--1.55 pp against 0.47). The gradualism model gave Table~\ref{tab:cbexp}. Because this specification was chosen after others failed, these results are exploratory.
\begin{table}[t]\centering\small\begin{tabular}{lrrrrr}\toprule
Bank & $n$ & Stated & Implied $\hat\tau$ & Mean inflation & Closer \\\midrule
GB & 59 & 2.0 & 2.51 & 2.20 & mean\\
CA & 92 & 2.0 & 1.89 & 1.82 & implied\\
NO & 67 & 2.5 & 2.68 & 1.97 & implied\\
SE & 60 & 2.0 & 1.36 & 1.20 & implied\\
AU & 104 & 2.5 & 2.63 & 2.46 & mean\\
MX & 68 & 3.0 & 3.27 & 4.11 & implied\\
CL & 76 & 3.0 & 3.96 & 3.17 & mean\\
IN & 34 & 4.0 & 4.39 & 4.84 & implied\\
PL & 64 & 2.5 & 1.63 & 2.13 & mean\\
\midrule
MAE & & & 0.45 & 0.47 & \\
\bottomrule\end{tabular}
\caption{Exploratory banks (gradualism model, realised inputs).}\label{tab:cbexp}\end{table}

\paragraph{Pre-registered confirmatory test.} Before compiling targets for new banks we fixed the sample rule, window rule, specification and success criterion (a document committed to the project repository). Five banks qualified (Table~\ref{tab:cbconf}). In the three banks whose average inflation deviated from target by more than 0.5 pp (New Zealand, Colombia, Turkey) the implied target was closer to the stated target in all three (MAE 0.50 against 2.34), meeting the criterion. In Israel and Hungary, where inflation averaged close to target, average inflation was much closer. Across both samples, the implied target was closer in all seven off-target cases.
\begin{table}[t]\centering\small\begin{tabular}{lrrrrr}\toprule
Bank & $n$ & Stated & Implied $\hat\tau$ & Mean inflation & Closer \\\midrule
IL & 47 & 2.0 & 3.49 & 1.77 & mean\\
NZ & 41 & 2.0 & 1.66 & 2.54 & implied\\
CO & 40 & 3.0 & 3.02 & 3.79 & implied\\
HU & 33 & 3.0 & 4.97 & 3.33 & mean\\
TR & 32 & 5.0 & 6.14 & 10.70 & implied\\
\midrule
MAE & & & 0.99 & 1.52 & \\
\bottomrule\end{tabular}
\caption{Confirmatory banks, pre-registered specification.}\label{tab:cbconf}\end{table}

\paragraph{Uncertainty.} Table~\ref{tab:boot} reports 90\% moving-block bootstrap intervals (blocks of 8 quarters, 20 replicates, a lighter search per replicate). Twelve of fourteen intervals contain the stated target. The width separates reliable from unreliable estimates: intervals are narrow for Australia, Canada, the United Kingdom, Colombia and Mexico, and very wide for Israel, Chile and Hungary; for India, Poland and Turkey some replicates drive the transmission parameter $\alpha$ towards zero, which makes $\tau=c/\alpha$ unbounded. Israel, the largest confirmatory miss, is flagged as unreliable by its interval. Mexico's interval excludes the stated 3\% narrowly, and Hungary's excludes it clearly.
\begin{table}[t]\centering\small\begin{tabular}{lrrrrcr}\toprule
Bank & Stated & 5\% & Median & 95\% & \hd{Covers\\stated} & \hd{Degenerate\\replicates} \\\midrule
AU & 2.5 & 2.06 & 2.63 & 3.03 & yes & 0\\
CA & 2.0 & 1.31 & 1.89 & 2.27 & yes & 1\\
CL & 3.0 & 0.99 & 3.63 & 5.32 & yes & 0\\
CO & 3.0 & 2.69 & 3.06 & 4.65 & yes & 0\\
GB & 2.0 & 1.79 & 2.51 & 2.77 & yes & 42\\
HU & 3.0 & 2.90 & 4.10 & 5.23 & yes & 0\\
IL & 2.0 & 0.43 & 1.45 & 4.43 & yes & 4\\
IN & 4.0 & 2.38 & 4.70 & 6.27 & yes & 29\\
MX & 3.0 & 3.18 & 3.47 & 3.89 & no & 0\\
NO & 2.5 & 1.79 & 2.68 & 3.54 & yes & 10\\
NZ & 2.0 & 1.65 & 2.84 & 3.45 & yes & 0\\
PL & 2.5 & 0.66 & 1.56 & 2.18 & no & 8\\
SE & 2.0 & 0.21 & 0.88 & 2.15 & yes & 0\\
TR & 5.0 & -29.88 & 7.46 & 10.17 & yes & 39\\
\bottomrule\end{tabular}
\caption{Bootstrap intervals for the implied inflation target.}\label{tab:boot}\end{table}

\section{Applications}\label{sec:apps}
For budget shares the decision moves the goal variable one-for-one ($\alpha=1$), there is one goal ($K=1$), and each actor-year is a decision. We report the revealed band $[\tau-b,\tau+b]$ and the step sizes when the share is above (step down) or below (step up) the band. A step that is not identified (no year strictly outside the band on that side) is shown as a dash; in that case the corresponding band edge is also only weakly determined (Proposition~\ref{prop:edge}).

\paragraph{Aid.} Official development assistance as a share of GNI (OECD DAC1, 1960--2024, net disbursement basis to 2017 and grant-equivalent from 2019; the 2018 change spans the break and is dropped), 23 bilateral donors with at least 30 annual decisions (Table~\ref{tab:aid}). Every donor has the same stated commitment, 0.7\% of GNI. Revealed targets differ widely: Norway (0.98), Luxembourg (1.00), Sweden (0.90) and Denmark (0.86) are above 0.7, the Netherlands is at 0.73, and the United States (0.17), Italy (0.14), Korea (0.14) and Spain (0.19) are far below. For the United Kingdom the upper edge of the revealed band is 0.70, the level written into law in 2015. Five donors have national ODA goals that we verified from OECD peer reviews and parliamentary documents after the estimates had been produced (Table~\ref{tab:aidgt}): Luxembourg set itself 1\% of GNI in 1999, Norway committed to 1\% in 2005, Sweden held a 1\% goal until its 2023 budget, Denmark's stated goal was 0.8\% in 2011, and the Netherlands set 0.7\% from 2012 after about 0.8\% before. In four of the five the revealed target is much closer to the national goal than the donor's average ODA level; for the Netherlands the two coincide. This comparison was not pre-registered, and the estimates use the whole period while the goals were adopted at different dates.
\begin{table}[t]\centering\small
\begin{tabular}{lcccc}\toprule
Donor & National goal (verified) & Revealed $\hat\tau$ & Mean level & Closer\\\midrule
Luxembourg & 1.0 (from 1999) & 1.00 & 0.63 & revealed\\
Norway & 1.0 (from 2005) & 0.98 & 0.80 & revealed\\
Sweden & 1.0 (until 2022) & 0.90 & 0.77 & revealed\\
Denmark & 0.8 (2011) & 0.86 & 0.70 & revealed\\
Netherlands & 0.7 from 2012 (0.8 before) & 0.73 & 0.73 & equal\\
\bottomrule\end{tabular}
\caption{Revealed aid targets against verified national goals (per cent of GNI). Sources: data/aid/targets.csv and docs/lit/14\_national\_oda\_targets.md in the replication files.}\label{tab:aidgt}\end{table}
\begin{table}[t]\centering\small\begin{tabular}{lrrrrrr}\toprule
Donor & $n$ & $\hat\tau$ & Band & Near-optimal $\tau$ set & Steps (down, up) & Mean level\\\midrule
LUX & 43 & 1.00 & [1.00, 1.00] & [1.00, 1.02] & 0.06, 0.03 & 0.63\\
NOR & 63 & 0.98 & [0.93, 1.03] & [0.95, 1.05] & 0.08, 0.03 & 0.80\\
SWE & 63 & 0.90 & [0.83, 0.97] & [0.87, 0.95] & 0.43, 0.06 & 0.77\\
DNK & 63 & 0.86 & [0.71, 1.01] & [0.77, 0.88] & 0.03, 0.02 & 0.70\\
FIN & 62 & 0.59 & [0.55, 0.64] & [0.48, 0.67] & 0.16, 0.01 & 0.34\\
DEU & 63 & 0.57 & [0.39, 0.76] & [0.36, 0.82] & 0.03, 0.01 & 0.41\\
FRA & 63 & 0.56 & [0.44, 0.69] & [0.54, 0.59] & 0.09, 0.01 & 0.55\\
TUR & 30 & 0.56 & [0.56, 0.56] & [0.20, 0.58] & 0.17, 0.01 & 0.25\\
NLD & 63 & 0.55 & [0.49, 0.61] & [0.49, 0.74] & 0.01, 0.07 & 0.73\\
CHE & 63 & 0.51 & [0.51, 0.51] & [0.29, 0.58] & 0.09, 0.01 & 0.30\\
AUS & 63 & 0.49 & [0.43, 0.56] & [0.25, 0.58] & 0.09, 0.00 & 0.39\\
BEL & 63 & 0.48 & [0.46, 0.50] & [0.42, 0.51] & 0.12, 0.04 & 0.48\\
IRL & 49 & 0.43 & [0.31, 0.55] & [0.35, 0.47] & 0.11, 0.01 & 0.31\\
CAN & 63 & 0.36 & [0.26, 0.47] & [0.35, 0.38] & 0.07, 0.02 & 0.35\\
AUT & 63 & 0.34 & [0.33, 0.35] & [0.25, 0.37] & 0.07, 0.02 & 0.23\\
NZL & 62 & 0.29 & [0.24, 0.34] & [0.28, 0.34] & 0.11, 0.01 & 0.26\\
JPN & 63 & 0.28 & [0.20, 0.37] & [0.25, 0.31] & 0.07, 0.06 & 0.25\\
GBR & 63 & 0.27 & [0.27, 0.27] & [0.04, 0.84] & 0.00, 0.03 & 0.43\\
PRT & 43 & 0.27 & [0.25, 0.28] & [0.25, 0.27] & 0.35, 0.01 & 0.21\\
USA & 63 & 0.21 & [0.19, 0.24] & [0.21, 0.22] & 0.03, 0.02 & 0.25\\
ESP & 43 & 0.18 & [0.10, 0.27] & [0.17, 0.28] & 0.03, 0.04 & 0.22\\
KOR & 36 & 0.15 & [0.15, 0.15] & [0.15, 0.18] & 0.00, 0.01 & 0.08\\
ITA & 63 & 0.14 & [0.13, 0.15] & [0.13, 0.28] & 0.02, 0.04 & 0.20\\
\bottomrule\end{tabular}
\caption{Revealed ODA/GNI aspiration bands by donor (per cent of GNI).}\label{tab:aid}\end{table}

\paragraph{Defence and leverage.} Table~\ref{tab:deffirms} pools allies (NATO definition, per cent of GDP) and, separately, US listed firms by two-digit industry (book leverage, one minus equity over assets; up to 3,000 firm-years per industry). For allies the only identified behaviour is raising spending when below the band: the lower edge is 1.74\% in 2015--2021 and 3.21\% in 2022--2025, with the step up rising from 0.03 to 0.12 points of GDP per year; the upper edges are not identified because almost no ally was above them. The shift of the lower edge is consistent with the move from the 2\% guideline towards the 3.5\% core target agreed in 2025, but the 2022 invasion of Ukraine changed the threat environment at the same time, and the two cannot be separated here. For firms the model is mostly uninformative: revealed bands span most of the leverage range and step sizes are not identified, except for depository institutions (SIC 60, band 0.87--1.00) and holding and investment offices (SIC 67). A one-goal model of book leverage does not capture firms' financing decisions. A two-goal version that adds the cash-to-assets ratio, with its transmission estimated, did not help either: in three large industries (SIC 28, 36, 73; 2,000 firm-years each) its mean absolute decision error was within 3\% of predicting no change at all (results/firms\_k2.csv).
\begin{table}[t]\centering\small\begin{tabular}{llrrrrr}\toprule
Domain & Unit & $n$ & Band & Step down & Step up & Mean level\\\midrule
Defence, NATO def. & 2015-2021 & 217 & [1.99, 3.48] & 0.20 & 0.03 & 1.48\\
Defence, NATO def. & 2022-2025 & 123 & [3.21, 3.66] & 0.21 & 0.12 & 1.85\\
Book leverage & SIC 13 & 1556 & [0.03, 0.83] & -- & -- & 0.47\\
Book leverage & SIC 28 & 3000 & [0.08, 0.82] & -- & -- & 0.37\\
Book leverage & SIC 35 & 1824 & [0.02, 0.93] & -- & 0.00 & 0.48\\
Book leverage & SIC 36 & 2930 & [0.02, 0.93] & -- & -- & 0.41\\
Book leverage & SIC 38 & 2812 & [0.07, 0.81] & -- & -- & 0.39\\
Book leverage & SIC 49 & 1543 & [0.03, 0.97] & -- & -- & 0.66\\
Book leverage & SIC 60 & 3000 & [0.87, 1.00] & -- & 0.01 & 0.89\\
Book leverage & SIC 63 & 1704 & [0.15, 0.98] & -- & -- & 0.75\\
Book leverage & SIC 67 & 3000 & [0.61, 1.23] & -- & 0.01 & 0.60\\
Book leverage & SIC 73 & 3000 & [0.08, 0.93] & 0.07 & 0.07 & 0.51\\
\bottomrule\end{tabular}
\caption{Defence spending (pooled allies) and firm leverage (pooled by industry).}\label{tab:deffirms}\end{table}

\section{Discussion and limitations}
The method recovers what actors pursue when outcomes and announcements diverge, which is where revealed aspirations matter most. It is less reliable than simple averages when outcomes already sit at the target. Inputs are proxies (realised inflation rather than real-time forecasts; official forecasts are conditional on policy paths and drift towards the target by construction). The forward model is myopic and linear in transmission. Estimation is not certified globally optimal for long series. The confirmatory sample is small. Revealed aspiration levels are those of a goal program that rationalises the decisions; they are not a claim about the decision maker's reasoning.

\bibliographystyle{plainnat}
\bibliography{refs}
\end{document}
